% TOP_PROBLEM: {"id": 30006748, "title": "Torsion Conjecture for Abelian Varieties of Dimension at Least Two", "category_id": 1, "category": {"id": 1, "name": "number_theory", "display_name": "Number Theory", "description": "Properties of integers, prime numbers, Diophantine equations.", "slug": "number-theory", "order_index": 1, "created_at": "2026-07-31T15:26:25.670Z"}, "status": "open"}
% FIRST_POSTED: null
% !TeX program = lualatex
% Compile twice: lualatex 30006748.tex
% All references and prose are embedded; no BibTeX database or external inputs.
\documentclass[11pt,a4paper]{article}
\usepackage[margin=24mm,headheight=14pt]{geometry}
\usepackage{fontspec}
\setmainfont{Latin Modern Roman}
\setsansfont{Latin Modern Sans}
\setmonofont{Latin Modern Mono}
\usepackage{amsmath,amssymb,mathtools}
\usepackage{unicode-math}
\setmathfont{Latin Modern Math}
\usepackage{microtype}
\usepackage{enumitem,xurl,needspace}
\usepackage[dvipsnames]{xcolor}
\usepackage{hyperref}
\hypersetup{unicode=true,colorlinks=true,linkcolor=MidnightBlue,urlcolor=MidnightBlue,
 pdftitle={Research attempt: Problem 30006748},pdfauthor={Research notebook},bookmarksnumbered=true}
\usepackage{bookmark}
\usepackage{tocloft}
\setlength{\cftsecnumwidth}{3em}
\cftsetpnumwidth{2.5em}
\cftsetrmarg{4em}
\usepackage{titlesec}
\titleformat{\section}{\Large\bfseries\raggedright}{\thesection}{0.7em}{}
\usepackage{fancyhdr}
\pagestyle{fancy}\fancyhf{}
\fancyhead[L]{\small\sffamily Open Problems Math | Research notebook}
\fancyhead[R]{\small Problem 30006748 | 27 September 2026}
\fancyfoot[C]{\thepage}
\setlength{\parindent}{0pt}
\setlength{\parskip}{3pt plus 1pt}
\setlength{\emergencystretch}{3em}
\setcounter{tocdepth}{1}
\setcounter{secnumdepth}{1}
\setlist[itemize]{leftmargin=3em,itemsep=5pt,topsep=4pt,parsep=0pt}
\allowdisplaybreaks
\newcommand{\N}{\mathbb{N}}
\newcommand{\Z}{\mathbb{Z}}
\newcommand{\Q}{\mathbb{Q}}
\newcommand{\R}{\mathbb{R}}
\newcommand{\C}{\mathbb{C}}
\newcommand{\F}{\mathbb{F}}
\newcommand{\A}{\mathbb{A}}
\newcommand{\PP}{\mathbb P}
\newcommand{\E}{\mathbb{E}}
\newcommand{\eps}{\varepsilon}
\newcommand{\dd}{\,\mathrm d}
\newcommand{\sref}[2]{\hyperlink{source:#1:#2}{[S#2]}}
\newcommand{\eref}[2]{\hyperlink{extra:#1:#2}{[E#2]}}
% Turn the next switch off for a shorter reading copy. The quotations preserve
% every exactTarget field, including qualifications omitted from a short summary.
\newif\ifshowcatalogue
\showcataloguetrue
\newcommand{\cataloguescope}[1]{\ifshowcatalogue
 \paragraph{Exact scope in the supplied catalogue.}
 \begin{quote}\small #1\end{quote}\fi}
\usepackage{etoolbox}
\AtBeginEnvironment{itemize}{\small}
\numberwithin{equation}{section}
% Standalone export. Original research content preserved.
\begin{document}
\setcounter{section}{0}
% ================================================================
% problemId: 30006748
% Canonical problem ID: 30006748
% exactTarget SHA-256: f8a7f42d292ab652354f834d31f85f0d05bd70426f1fb49b0ffd4a1a2a86fb4d
\clearpage
\section*{\texorpdfstring{Torsion Conjecture for Abelian Varieties of Dimension at Least Two}{Torsion Conjecture for Abelian Varieties of Dimension at Least Two}}\label{30006748}
\subsection{Definitions and mathematical statement}
For an abelian variety $A/K$, its rational torsion subgroup is
\[
 A(K)_{\rm tors}=\{P\in A(K):mP=0\text{ for some integer }m\ge1\}.
\]
For each fixed number field $K$ and dimension $g\ge2$, the conjecture asks for a finite $B(g,K)$ with $\#A(K)_{\rm tors}\le B(g,K)$ for every $g$-dimensional $A/K$. The bound cannot depend on height, endomorphism ring, polarization, or reduction type. It need not be effective, and uniformity over all fields of a fixed degree is not appended.
\par\smallskip\textit{Formulation and concept references:} \sref{30006748}{1}.
\cataloguescope{For every integer g>=2 and every number field K, there exists a finite constant B(g,K) such that every abelian variety A defined over K of dimension g satisfies \#A(K)\_tors <= B(g,K). Here A(K)\_tors is the subgroup of K-rational torsion points. The bound is uniform over all such A, without dependence on A, its height, endomorphisms, polarization degree or reduction type; it may depend on g and K. No effective formula, explicit numerical value or classification of torsion groups is required.}
\subsection{Short English statement}
Is rational torsion uniformly bounded across all abelian varieties of a fixed dimension over a fixed number field?
% BEGIN RESEARCH ADDITION 141
\subsection{Research attempt}

\paragraph{Current frontier.}
The number-field and function-field claims of the supplied 2026 manuscript
must be separated. Its Theorem 1.3 concerns characteristic-zero function
fields with a trace condition; Theorem 1.8 retains reduction/height quantities,
and Theorem 1.9 assumes its generalized Szpiro Conjecture 1.7 for the
number-field conclusion in the semistable case. None is an unconditional
proof for all $A/K$. These are the manuscript's stated results, not an
independent verification of its proofs. \sref{30006748}{1}
A 2026 geometric approach to $\ell$-primary torsion concerns fibres of an
abelian scheme over a curve, not all primes and all varieties at once.
\eref{30006748}{2}

\paragraph{Attack route.}
Combine several good reductions while retaining prime-primary information.
Uniformly small good places would yield a bound without height estimates.
We derive an exact certificate and then disprove that proposed place-selection
lemma. Two alternatives are to pass to a principally polarized variety or
through an isogeny; their costs are different and must be tracked.

\paragraph{Attempt: a multi-place divisibility certificate.}
Let $T=A(K)_{\rm tors}$ and let $V$ be a finite set of good places with at
least two distinct residue characteristics. Write $p_v=\operatorname{char}k_v$,
$q_v=\#k_v$ and $a_v=\#A(k_v)$. Then
\begin{equation}\label{30006748:certificate}
 |T|\mid D_V:=\prod_{\ell\ {\rm prime}}
       \ell^{\min_{v\in V,\ p_v\ne\ell}v_\ell(a_v)}.
\end{equation}
Each minimum is nonempty, and only finitely many exponents are nonzero.
At a good place the reduction kernel is a formal group with additive
residue-field successive quotients, so it has no prime-to-$p_v$ torsion.
Thus $T[\ell^\infty]$ injects into $A(k_v)$ if $\ell\ne p_v$; taking
valuations proves the formula. The local and finite-field ingredients are
reviewed in Clark--Xarles, Sections 3.1--3.2. \eref{30006748}{1}

For two places of characteristics $p\ne q$, with point counts $a,b$,
\begin{equation}\label{30006748:two}
 D_{v,w}=\gcd(a,b)\,
 p^{\max(0,v_p(b)-v_p(a))}\,
 q^{\max(0,v_q(a)-v_q(b))}.
\end{equation}
In particular
\[
 |T|\le ab\le\bigl((1+\sqrt{q_v})(1+\sqrt{q_w})\bigr)^{2g}.
\]
The last step is the Weil point bound at good reduction, not a bound on a
bad special fibre. \eref{30006748}{1} The residue-characteristic corrections in
\eqref{30006748:two} cannot simply be dropped by taking a greatest common divisor.
For $E:y^2=x(x-1)(x-2)$ and $A=E^g$, the companion program computes odd-prime
point counts, applies the certificate, and checks divisibility by $4^g$
from the visible rational two-torsion subgroup. It is an exact finite
certificate computation, not an assertion of exact torsion for every input.

\paragraph{A small-prime escape construction.}
The proposed uniform lemma ``two good places of different characteristics
have norm at most $Q(g,K)$'' fails already over $\Q$. Given $B\ge3$, take
\[
 t=2\prod_{\substack{3\le p\le B\\p\ {\rm prime}}}p,
 \qquad E_t:y^2=x(x-1)(x-t).
\]
Its invariants are
\[
 c_4=16(t^2-t+1),\qquad \Delta=16t^2(t-1)^2.
\]
For every odd $p\le B$, $v_p(c_4)=0$ and $v_p(\Delta)>0$. The nonintegral
$j=c_4^3/\Delta$ excludes good reduction. The integral equation is minimal
there since a nonminimal scaling would require $v_p(c_4)\ge4$; its reduction
is nodal. Thus $E_t^g$ has bad reduction at every such prime, also evident
from its nontrivial direct-sum inertia action. Even if $2$ were good, it
would not supply two different small residue characteristics. No large
rational torsion is asserted for these examples.

There is a local obstruction to weakening good reduction to semistability.
For a Tate curve over $\Q_p$ with parameter $q=p^m$, uniformization gives
$E_q(\Q_p)=\Q_p^\times/q^{\Z}$. The point represented by $p$ has order exactly
$m$, because $p^r\in p^{m\Z}$ precisely when $m\mid r$. Arbitrary
semistable curves over a fixed local field therefore have no uniform torsion
bound. The split-multiplicative distinction is essential in the local
results. \eref{30006748}{1} These are local points, not rational points over a
fixed number field, and do not contradict the catalogue target.

\paragraph{Two transfers with different costs.}
Zarhin's construction $Z(A)=(A\times A^\vee)^4$ has dimension $8g$ and a
principal polarization over $K$; the line-bundle construction is in the
source's Appendix B. \sref{30006748}{1} Since $A$ embeds as a factor over $K$,
a uniform bound for principally polarized varieties in dimension $8g$
would imply the target in dimension $g$, without a degree-of-polarization
loss. That change of dimension is allowed by the quantifiers.

In contrast, for a $K$-isogeny $\varphi:A\to B$ of degree $d$,
\[
 |A(K)_{\rm tors}|\le d\,|B(K)_{\rm tors}|.
\]
The kernel has at most $d$ rational points and the image of torsion is
torsion. Decomposition up to isogeny into elliptic factors is therefore not
a uniform argument unless the isogeny degree is controlled. The factor
embedding in Zarhin's construction avoids this particular loss.

\paragraph{Stress test.}
The certificate handles all primary parts; it does not assume injectivity
on the residue-characteristic primary subgroup. Taking more places gives
curve-dependent improvements but no uniform theorem without a uniform
selection mechanism. The Legendre and Tate examples disprove two specific
local shortcuts, not the global assertion. Bounds for each fixed $\ell$
would still require a uniform exclusion of sufficiently large torsion
primes before they could bound the whole torsion group. No such exclusion
has been obtained.

\paragraph{Outcome.}
\textbf{Outcome: Exploratory approach with unresolved gap.}
The attempt proves an exact finite-reduction certificate and explains a
polarization-independent transfer, together with explicit failures of two
local strategies. These elementary deductions are not claimed to be new
in the literature. A genuinely global uniform constraint remains necessary
when all small places have bad reduction; the conjectural height/reduction
estimates in the source are not proved here.

% END RESEARCH ADDITION 141
\subsection{Sources}
\begin{itemize}\raggedright
\item[\textnormal{[S1]}]\hypertarget{source:30006748:1}{} Uniformity in rational torsion and small points on abelian varieties; arXiv2608.22285v2.
\url{https://arxiv.org/pdf/2608.22285v2}.
{\small\textit{Catalogue formulation source.}}
% BEGIN RESEARCH REFERENCES 141
\item[\textnormal{[E1]}]\hypertarget{extra:30006748:1}{} P. L. Clark and X. Xarles,
\emph{Local bounds for torsion points on abelian varieties}, Canadian Journal
of Mathematics 60 (2008), 532--555, Introduction, Sections 2.3, 3.1,
Proposition 3.3; DOI 10.4153/CJM-2008-026-x.
\url{https://www.cambridge.org/core/services/aop-cambridge-core/content/view/FE99F6239E5F5B8995C4DEDB1B606EB9/S0008414X00019647a.pdf/local_bounds_for_torsion_points_on_abelian_varieties.pdf}.
\item[\textnormal{[E2]}]\hypertarget{extra:30006748:2}{} Z. Ji, J. Song and J. Xie,
\emph{A geometric approach to the uniform boundedness of $\ell$-primary
torsion points}, arXiv:2601.15089, 21 January 2026; abstract and application
to abelian schemes over curves. \url{https://arxiv.org/abs/2601.15089}.

% END RESEARCH REFERENCES 141
\end{itemize}

\end{document}
